---
tipo: Discursiva
dificuldade: Média
disciplina: Matemática
assunto: Semelhança de Triângulos
tags: [bissetriz, segmentos proporcionais, triângulo]
gabarito:
fonte: material didático
---
\question
No triângulo ABC da figura abaixo, AS é bissetriz interna do ângulo  e AP é bissetriz externa. Calcule a medida do segmento SP.

% figura presente no enunciado

---
tipo: Discursiva
dificuldade: Média
disciplina: Matemática
assunto: Razões e Proporções
tags: [escala, mapa, proporcionalidade]
gabarito:
fonte: material didático
---
\question
A figura abaixo mostra um fragmento de mapa, em que se vê o trecho reto da estrada que liga as cidades de Paraguaçu e Piripiri. Os números apresentados no mapa representam as distâncias, em quilômetros, entre cada cidade e o ponto de início da estrada (que não aparece na figura). Os traços perpendiculares à estrada estão uniformemente espaçados de 1 cm.
\begin{parts}
\part Para representar a escala de um mapa, usamos a notação 1 : X, onde X é a distância real correspondente à distância de 1 unidade do mapa. Usando essa notação, indique a escala do mapa dado acima.
\part Repare que há um posto exatamente sobre um traço perpendicular à estrada. Em que quilômetro (medido a partir do ponto de início da estrada) encontra-se tal posto?
\part Imagine que você tenha que reproduzir o mapa dado usando a escala 1 : 500.000. Se você fizer a figura em uma folha de papel, qual será a distância, em centímetros, entre as cidades de Paraguaçu e Piripiri?
\end{parts}

% figura presente no enunciado

---
tipo: Múltipla Escolha
dificuldade: Média
disciplina: Matemática
assunto: Semelhança de Triângulos
tags: [bissetriz, triângulo retângulo, proporcionalidade]
gabarito:
fonte: material didático
---
\question
Na figura, o triângulo ABD é reto em B, e AC é a bissetriz de BAD. Se AB = 2BC, fazendo BC = b e CD = d, então:
\begin{choices}
\choice $d = b$
\choice $d = \frac{5}{2}b$
\choice $d = \frac{5}{3}b$
\choice $d = \frac{6}{5}b$
\choice $d = \frac{5}{4}b$
\end{choices}

% figura presente no enunciado

---
tipo: Discursiva
dificuldade: Média
disciplina: Matemática
assunto: Teorema de Tales
tags: [feixe de paralelas, terrenos, proporcionalidade]
gabarito:
fonte: material didático
---
\question
Três terrenos têm frente para a rua “A” e para a rua “B”, como na figura. As divisas laterais são perpendiculares à rua “A”. Qual a medida de frente para a rua “B” de cada lote, sabendo que a frente total para essa rua é 180 m?

% figura presente no enunciado

---
tipo: Discursiva
dificuldade: Fácil
disciplina: Matemática
assunto: Semelhança de Triângulos
tags: [razão de semelhança, perímetro, triângulos semelhantes]
gabarito:
fonte: material didático
---
\question
Um triângulo têm lados medindo 4 cm, 5 cm e 7 cm. Um segundo triângulo, semelhante ao primeiro, tem perímetro 64 cm. Determine as medidas dos lados do segundo triângulo.

---
tipo: Discursiva
dificuldade: Média
disciplina: Matemática
assunto: Área das Figuras Planas
tags: [triângulo isósceles, círculo inscrito, raio]
gabarito:
fonte: material didático
---
\question
Um triângulo isósceles tem lados medindo 10 cm, 10 cm e 12 cm. A altura relativa ao maior lado mede 8 cm. Ache o raio do círculo inscrito.

---
tipo: Múltipla Escolha
dificuldade: Média
disciplina: Matemática
assunto: Semelhança de Triângulos
tags: [rampa, proporcionalidade, altura]
gabarito:
fonte: material didático
---
\question
A rampa de um hospital tem na sua parte mais elevada uma altura de 2,2 metros. Um paciente ao caminhar sobre a rampa percebe que se deslocou 3,2 metros e alcançou uma altura de 0,8 metro.
A distância em metros que o paciente ainda deve caminhar para atingir o ponto mais alto da rampa é
\begin{choices}
\choice 1,16 metro.
\choice 3,0 metros.
\choice 5,4 metros.
\choice 5,6 metros.
\choice 7,04 metros.
\end{choices}

---
tipo: Discursiva
dificuldade: Difícil
disciplina: Matemática
assunto: Semelhança de Triângulos
tags: [círculos tangentes, triângulo isósceles, altura]
gabarito:
fonte: material didático
---
\question
Dois círculos tangentes exteriormente, de raios 2R e R, estão inscritos num triângulo isósceles ABC, conforme a figura a seguir. Determinar a altura AH.

% figura presente no enunciado

---
tipo: Discursiva
dificuldade: Média
disciplina: Matemática
assunto: Semelhança de Triângulos
tags: [retângulos, triângulo retângulo, proporcionalidade]
gabarito:
fonte: material didático
---
\question
No interior de um triângulo retângulo foram colocados dois retângulos iguais, como mostra a figura abaixo. Se cada retângulo tem dimensões, em cm, de 6 e 16, calcule o valor de x.

% figura presente no enunciado

---
tipo: Múltipla Escolha
dificuldade: Média
disciplina: Matemática
assunto: Teorema de Tales
tags: [paralelogramo, segmentos proporcionais, paralelas]
gabarito:
fonte: material didático
---
\question
Considere o paralelogramo MNPQ representado na figura abaixo.
Se $QT = \frac{1}{3}QP$, então:
\begin{choices}
\choice $MV = \frac{1}{2}VP$
\choice $MV = 3VP$
\choice $MV = 2VP$
\choice $MV = \frac{2}{9}VP$
\choice $MV = \frac{3}{2}VP$
\end{choices}

% figura presente no enunciado

---
tipo: Discursiva
dificuldade: Média
disciplina: Matemática
assunto: Semelhança de Triângulos
tags: [altura relativa, triângulo retângulo, segmentos]
gabarito:
fonte: material didático
---
\question
O triângulo ABC da figura a seguir tem ângulo reto em B. O segmento BD é a altura relativa a AC. Os segmentos AD e DC medem 12 cm e 4 cm, respectivamente. O ponto E pertence ao lado BC e BC = 4EC.
Determine o comprimento do segmento DE.

% figura presente no enunciado

---
tipo: Múltipla Escolha
dificuldade: Média
disciplina: Matemática
assunto: Área das Figuras Planas
tags: [retângulo, diagonal, área de triângulo]
gabarito:
fonte: material didático
---
\question
A figura representa um retângulo ABCD, com AB = 5 e AD = 3. O ponto E está no segmento CD de maneira que CE = 1, e F é o ponto de interseção da diagonal AC com o segmento BE. Então, a área do triângulo BCF vale
\begin{choices}
\choice $\frac{6}{5}$
\choice $\frac{5}{4}$
\choice $\frac{4}{3}$
\choice $\frac{7}{5}$
\choice $\frac{3}{2}$
\end{choices}

% figura presente no enunciado

---
tipo: Múltipla Escolha
dificuldade: Média
disciplina: Matemática
assunto: Semelhança de Triângulos
tags: [baricentro, paralelas, perímetro]
gabarito:
fonte: material didático
---
\question
Considere o triângulo ABC, de lados AB = 15, AC = 10, BC = 12 e seu baricentro G. Traçam-se GE e GF paralelos a AB e AC, respectivamente, conforme a figura abaixo. O perímetro do triângulo GEF é um número que, escrito na forma de fração irredutível, tem a soma do numerador com o denominador igual a:
\begin{choices}
\choice 43
\choice 38
\choice 40
\choice 35
\end{choices}

% figura presente no enunciado

---
tipo: Múltipla Escolha
dificuldade: Média
disciplina: Matemática
assunto: Semelhança de Triângulos
tags: [triângulos retângulos, segmentos proporcionais, medida]
gabarito:
fonte: material didático
---
\question
Na figura ABC e CDE são triângulos retângulos, AB = 1, BC = 3 e BE = 2DE. Logo, a medida de AE é:
\begin{choices}
\choice $\frac{3}{2}$
\choice $\frac{5}{2}$
\choice $\frac{7}{2}$
\choice $\frac{11}{2}$
\choice $\frac{13}{2}$
\end{choices}

% figura presente no enunciado

---
tipo: Discursiva
dificuldade: Média
disciplina: Matemática
assunto: Área das Figuras Planas
tags: [retângulo inscrito, triângulo, perímetro]
gabarito:
fonte: material didático
---
\question
Na figura, o triângulo ABC tem base BC = 20 cm e altura 12 cm. Se DE // BC, e se o retângulo DEFG tem perímetro 30 cm, ache os lados do retângulo.

% figura presente no enunciado

---
tipo: Discursiva
dificuldade: Média
disciplina: Matemática
assunto: Semelhança de Triângulos
tags: [sombra, ladeira, poste]
gabarito:
fonte: material didático
---
\question
Um homem, de 1,80 m de altura, sobe uma ladeira com inclinação de 30°, conforme mostra a figura. No ponto A está um poste vertical de 5 metros de altura, com uma lâmpada no ponto B. Pede-se para:
\begin{parts}
\part Calcular o comprimento da sombra do homem depois que ele subiu 4 metros ladeira acima.
\part Calcular a área do triângulo ABC.
\end{parts}

% figura presente no enunciado

---
tipo: Discursiva
dificuldade: Média
disciplina: Matemática
assunto: Semelhança de Triângulos
tags: [poste, sombra, edifício]
gabarito:
fonte: material didático
---
\question
Um poste tem uma lâmpada colocada a 4 m de altura. Um homem de 2 m de altura caminha, a partir do poste em linha reta, em direção à porta de um edifício que está a uma distância de 28 m do poste. Calcule o comprimento da sombra do homem que é projetada sobre a porta do edifício, no instante em que ele está a 10,5 m dessa porta. Sua resposta deve vir acompanhada de um desenho ilustrativo da situação descrita.

---
tipo: Discursiva
dificuldade: Difícil
disciplina: Matemática
assunto: Circunferência e Círculo
tags: [esferas, centros, caminho mais curto]
gabarito:
fonte: material didático
---
\question
Três goiabas perfeitamente esféricas de centros $C_1$, $C_2$ e $C_3$ e raios 2 cm, 8 cm e 2 cm estão sobre uma mesa tangenciando-se como sugere a figura abaixo.
Um bichinho que está no centro da primeira goiaba quer se dirigir para o centro da terceira pelo caminho mais curto. Quantos centímetros percorrerá?

% figura presente no enunciado

---
tipo: Múltipla Escolha
dificuldade: Média
disciplina: Matemática
assunto: Área das Figuras Planas
tags: [retângulo inscrito, base e altura, fórmula]
gabarito:
fonte: material didático
---
\question
O triângulo ABC tem altura h e base b (ver figura). Nele, está inscrito o retângulo DEFG, cuja base é o dobro da altura. Nessas condições, a altura do retângulo, em função de h e b, é dada pela fórmula:
\begin{choices}
\choice $\frac{bh}{b+h}$
\choice $\frac{2bh}{b+h}$
\choice $\frac{bh}{2b+h}$
\choice $\frac{bh}{b+2h}$
\choice $\frac{bh}{2(h+b)}$
\end{choices}

% figura presente no enunciado

---
tipo: Múltipla Escolha
dificuldade: Média
disciplina: Matemática
assunto: Teorema de Tales
tags: [circuito triangular, paralelas, perímetro]
gabarito:
fonte: material didático
---
\question
O circuito triangular de uma corrida está esquematizado na figura a seguir:
As ruas TP e SQ são paralelas. Partindo de S, cada corredor deve percorrer o circuito passando, sucessivamente, por R, Q, P e T, retornando, finalmente, a S.
Assinale a opção que indica o perímetro do circuito.
\begin{choices}
\choice 4,5 km
\choice 19,5 km
\choice 20 km
\choice 22,5 km
\choice 24,0 km
\end{choices}

% figura presente no enunciado

---
tipo: Discursiva
dificuldade: Difícil
disciplina: Matemática
assunto: Circunferência e Círculo
tags: [circunferências tangentes, raios, reta tangente]
gabarito:
fonte: material didático
---
\question
As quatro circunferências da figura são tangentes duas a duas e tangentes a reta r. Sabendo-se que os raios das duas menores medem 1 cm e 5 cm, determine o raio da maior.

% figura presente no enunciado

---
tipo: Múltipla Escolha
dificuldade: Difícil
disciplina: Matemática
assunto: Circunferência e Círculo
tags: [circunferências inscritas, triângulo equilátero, raio]
gabarito:
fonte: material didático
---
\question
Na figura a seguir, as circunferências têm raios iguais a R e estão inscritas em um triângulo equilátero de lado igual a 2 cm.
Assinale a alternativa que representa o valor de R.
\begin{choices}
\choice $R = \frac{1}{1+\sqrt{3}}$ cm
\choice $R = \frac{3}{1+\sqrt{3}}$ cm
\choice $R = \frac{3}{1+\sqrt{3}}$ cm
\choice $R = \frac{3}{2+\sqrt{3}}$ cm
\choice $R = \frac{2}{2+\sqrt{3}}$ cm
\end{choices}

% figura presente no enunciado

---
tipo: Múltipla Escolha
dificuldade: Média
disciplina: Matemática
assunto: Circunferência e Círculo
tags: [semicircunferência, triângulo isósceles, tangência]
gabarito:
fonte: material didático
---
\question
No triângulo isósceles PQR, da figura abaixo, RH é a altura relativa ao lado PQ.
Se M é o ponto médio de PR, então a semicircunferência de centro M, e tangente a RH em T tem raio r igual a:
\begin{choices}
\choice 0,50 cm
\choice 0,75 cm
\choice 0,90 cm
\choice 1,00 cm
\choice 1,50 cm
\end{choices}

% figura presente no enunciado

---
tipo: Múltipla Escolha
dificuldade: Média
disciplina: Matemática
assunto: Semelhança de Triângulos
tags: [futebol, trajetória retilínea, distância mínima]
gabarito:
fonte: material didático
---
\question
Um lateral L faz um lançamento para um atacante A, situado 32 m à sua frente em uma linha paralela à lateral do campo de futebol. A bola, entretanto, segue uma trajetória retilínea, mas não paralela à lateral e quando passa pela linha de meio do campo está a uma distância de 12 m da linha que une o lateral ao atacante. Sabendo-se que a linha de meio do campo está à mesma distância dos dois jogadores, a distância mínima que o atacante terá que percorrer para encontrar a trajetória da bola será de:
\begin{choices}
\choice 18,8 m
\choice 19,2 m
\choice 19,6 m
\choice 20 m
\choice 20,4 m
\end{choices}

% figura presente no enunciado

---
tipo: Discursiva
dificuldade: Média
disciplina: Matemática
assunto: Semelhança de Triângulos
tags: [rampa, inclinação constante, altura]
gabarito:
fonte: material didático
---
\question
Uma rampa de inclinação constante, como a que dá acesso ao Palácio do Planalto em Brasília, tem 4 metros de altura na sua parte mais alta. Uma pessoa, tendo começado a subi-la, nota que após caminhar 12,3 metros sobre a rampa está a 1,5 metro de altura em relação ao solo.
\begin{parts}
\part Faça uma figura ilustrativa da situação descrita.
\part Calcule quantos metros a pessoa ainda deve caminhar para atingir o ponto mais alto da rampa.
\end{parts}

---
tipo: Discursiva
dificuldade: Difícil
disciplina: Matemática
assunto: Circunferência e Círculo
tags: [corda, círculos tangentes, reta tangente]
gabarito:
fonte: material didático
---
\question
Sejam $C_1$, $C_2$ e $C_3$ três círculos de mesmo raio R e cujos centros $O_1$, $O_2$ e $O_3$ estão sobre uma mesma reta. Além disso, $C_1$ é tangente a $C_2$ e $C_2$ é tangente a $C_3$. Considere a reta D que passa por A e é tangente ao círculo $C_3$ (ver figura). Expresse o comprimento da corda BC, determinada por D em $C_2$, em função de R.

% figura presente no enunciado

---
tipo: Discursiva
dificuldade: Média
disciplina: Matemática
assunto: Semelhança de Triângulos
tags: [obelisco, sombra, proporcionalidade]
gabarito:
fonte: material didático
---
\question
Um obelisco de 12 m de altura projeta, num certo momento, uma sombra de 4,8 m de extensão. Calcule a distância máxima que uma pessoa de 1,80 m de altura poderá se afastar do centro da base do obelisco, ao longo da sombra, para, em pé, continuar totalmente na sombra.

---
tipo: Discursiva
dificuldade: Fácil
disciplina: Matemática
assunto: Semelhança de Triângulos
tags: [torre, sombras, proporcionalidade]
gabarito:
fonte: material didático
---
\question
Um menino de 1,50 m de altura observa, num dia de sol, as sombras de uma torre de rádio emissora e a sua própria. Não dispondo de fita metálica ou de trena, ele toma um cordão, mede sua sombra e a compara com a da torre, verificando que esta está 10 vezes maior que a sua. Calcule a altura da torre.

---
tipo: Discursiva
dificuldade: Média
disciplina: Matemática
assunto: Semelhança de Triângulos
tags: [eclipse solar, raio do sol, proporcionalidade]
gabarito:
fonte: material didático
---
\question
Num eclipse do Sol, o disco lunar cobre exatamente o disco solar, o que comprova que o ângulo sob o qual vemos o Sol é o mesmo sob o qual vemos a Lua. Considerando que o raio da lua é de 1.738 km e que a distância da Lua ao Sol é 400 vezes da terra à lua, calcule o raio do Sol.

---
tipo: Discursiva
dificuldade: Média
disciplina: Matemática
assunto: Semelhança de Triângulos
tags: [bissetriz interna, perímetro, lados]
gabarito:
fonte: material didático
---
\question
O perímetro de um triângulo ABC é 100 cm. Sabendo que a bissetriz do ângulo interno A divide o lado oposto BC em dois segmentos de 13,5 cm e 22,5 cm, determine as medidas dos lados desse triângulo.

---
tipo: Discursiva
dificuldade: Fácil
disciplina: Matemática
assunto: Semelhança de Triângulos
tags: [perímetro, triângulos semelhantes, razão]
gabarito:
fonte: material didático
---
\question
Um triângulo, cujos lados medem 12 m, 18 m e 20 m, é semelhante a outro cujo perímetro mede 10 m. Calcule a medida dos lados do triângulo menor.

---
tipo: Múltipla Escolha
dificuldade: Média
disciplina: Matemática
assunto: Teorema de Tales
tags: [quarteirões, ruas, proporcionalidade]
gabarito:
fonte: material didático
---
\question
No desenho acima apresentado, as frentes para a rua A dos quarteirões I e II medem, respectivamente, 250 m e 200 m, e a frente do quarteirão I para a rua B mede 40m a mais do que a frente do quarteirão II para a mesma rua. Sendo assim, pode-se afirmar que a medida, em metros, da frente do menor dos dois quarteirões para a rua B é:
\begin{choices}
\choice 160
\choice 180
\choice 200
\choice 220
\choice 240
\end{choices}

% figura presente no enunciado

---
tipo: Discursiva
dificuldade: Média
disciplina: Matemática
assunto: Teorema de Tales
tags: [retas paralelas, segmentos, incógnitas]
gabarito:
fonte: material didático
---
\question
Determine x e y, sendo r, s e t retas paralelas.

% figura presente no enunciado

---
tipo: Discursiva
dificuldade: Média
disciplina: Matemática
assunto: Semelhança de Triângulos
tags: [bissetriz, perímetro, lado]
gabarito:
fonte: material didático
---
\question
Determine a medida ao lado AB do triângulo ABC:
\begin{parts}
\part AS é bissetriz e o perímetro do $\triangle ABC$ é 75 cm.
\part AP é bissetriz do ângulo externo em A e o perímetro do $\triangle ABC$ é 23 m.
\end{parts}

% figura presente no enunciado

---
tipo: Discursiva
dificuldade: Média
disciplina: Matemática
assunto: Teorema de Tales
tags: [retas paralelas, segmentos proporcionais, incógnitas]
gabarito:
fonte: material didático
---
\question
Determine x e y nos casos:
\begin{parts}
\part 
\part 
\end{parts}

% figura presente no enunciado

---
tipo: Discursiva
dificuldade: Média
disciplina: Matemática
assunto: Teorema de Tales
tags: [retas paralelas, segmentos, incógnita]
gabarito:
fonte: material didático
---
\question
Sendo r e s retas paralelas, determine x nos casos:
\begin{parts}
\part
\part
\end{parts}

% figura presente no enunciado

---
tipo: Discursiva
dificuldade: Média
disciplina: Matemática
assunto: Semelhança de Triângulos
tags: [bissetriz, segmentos proporcionais, incógnita]
gabarito:
fonte: material didático
---
\question
Na figura abaixo, determine o valor de x.

% figura presente no enunciado

---
tipo: Discursiva
dificuldade: Média
disciplina: Matemática
assunto: Teorema de Tales
tags: [paralelogramo, transversal, segmentos]
gabarito:
fonte: material didático
---
\question
Calcule o valor de x, sabendo que a figura abaixo é um paralelogramo.

% figura presente no enunciado

---
tipo: Discursiva
dificuldade: Média
disciplina: Matemática
assunto: Semelhança de Triângulos
tags: [segmentos, triângulos, incógnita]
gabarito:
fonte: material didático
---
\question
Dada a figura, determine o valor de x.

% figura presente no enunciado

---
tipo: Discursiva
dificuldade: Média
disciplina: Matemática
assunto: Área das Figuras Planas
tags: [quadrados, lado, incógnita]
gabarito:
fonte: material didático
---
\question
Na figura abaixo, consideremos os quadrados de lados a e b $(a > b)$. Calcule o valor de x.

% figura presente no enunciado

---
tipo: Discursiva
dificuldade: Média
disciplina: Matemática
assunto: Área das Figuras Planas
tags: [quadrados, perímetro, lado]
gabarito:
fonte: material didático
---
\question
Na figura abaixo, consideremos os quadrados de lados x, 6 e 9. Determine o perímetro do quadrado de lado x.

% figura presente no enunciado

---
tipo: Discursiva
dificuldade: Média
disciplina: Matemática
assunto: Área das Figuras Planas
tags: [quadrado, lado, segmentos]
gabarito:
fonte: material didático
---
\question
Determine a medida do lado do quadrado da figura abaixo:

% figura presente no enunciado

---
tipo: Múltipla Escolha
dificuldade: Média
disciplina: Matemática
assunto: Área das Figuras Planas
tags: [losango inscrito, triângulo, lado]
gabarito:
fonte: material didático
---
\question
O losango ADEF está inscrito no triângulo ABC, como mostra a figura. Se AB = 12 m, BC = 8 m e AC = 6 m, o lado do losango mede:
\begin{choices}
\choice 5m
\choice 3m
\choice 2m
\choice 4m
\choice 8m
\end{choices}

% figura presente no enunciado

---
tipo: Discursiva
dificuldade: Média
disciplina: Matemática
assunto: Circunferência e Círculo
tags: [círculo inscrito, triângulo, segmentos]
gabarito:
fonte: material didático
---
\question
Considerando as medidas indicadas na figura e sabendo que o círculo está inscrito no triângulo, determine x.

% figura presente no enunciado

---
tipo: Discursiva
dificuldade: Média
disciplina: Matemática
assunto: Semelhança de Triângulos
tags: [bissetriz interna, bissetriz externa, perímetro]
gabarito:
fonte: material didático
---
\question
Consideremos um triângulo ABC de 15 cm de perímetro. A bissetriz externa do ângulo  desse triângulo encontra o prolongamento do lado BC em um ponto S. Sabendo que a bissetriz interna do ângulo  determina sobre BC dois segmentos BP e PC de medidas 3 cm e 2 cm, respectivamente, determine as medidas dos lados do triângulo e a medida do segmento CS.

---
tipo: Discursiva
dificuldade: Média
disciplina: Matemática
assunto: Semelhança de Triângulos
tags: [trapézio, prolongamento dos lados, alturas]
gabarito:
fonte: material didático
---
\question
Prolongando-se os lados oblíquos às bases de um trapézio, obtemos um ponto E os triângulos ECD e EAB. Determine a relação entre as alturas dos dois triângulos, relativas aos lados que são bases do trapézio, sendo 12 cm e 4 cm as medidas das bases do trapézio.

---
tipo: Discursiva
dificuldade: Média
disciplina: Matemática
assunto: Semelhança de Triângulos
tags: [trapézio, triângulos semelhantes, alturas]
gabarito:
fonte: material didático
---
\question
As bases de um trapézio ABCD medem 50 cm e 30 cm e a altura 10 cm. Prolongando-se os lados não paralelos, eles se interceptam num ponto E. Determine a altura ABE do triângulo ABE e a altura do triângulo CDE (vide figura).

% figura presente no enunciado

---
tipo: Discursiva
dificuldade: Média
disciplina: Matemática
assunto: Semelhança de Triângulos
tags: [triângulo retângulo, segmentos, medida]
gabarito:
fonte: material didático
---
\question
Na figura, temos: AB = 8, BC = 15, AC = 17 e EC = 4. Determine x = DE e y = CD.

% figura presente no enunciado

---
tipo: Discursiva
dificuldade: Média
disciplina: Matemática
assunto: Área das Figuras Planas
tags: [retângulo, diagonal, ponto médio]
gabarito:
fonte: material didático
---
\question
Num retângulo ABCD, os lados AB e BC medem 20 cm e 12 cm, respectivamente. Sabendo que M é o ponto médio do lado AB, calcule EF, distância do ponto E ao lado AB, sendo E a interseção da diagonal BD com o segmento CM.

% figura presente no enunciado

---
tipo: Múltipla Escolha
dificuldade: Difícil
disciplina: Matemática
assunto: Semelhança de Triângulos
tags: [mediana, paralela, soma de segmentos]
gabarito:
fonte: material didático
---
\question
Considere o triângulo ABC, onde D é a mediana relativa ao lado BC. Por um ponto arbitrário M do segmento BD, tracemos o segmento MP paralelo a AD, onde P é o ponto de interseção desta paralela com o prolongamento do lado AC (figura). Se N é o ponto de interseção de AB com MP, podemos afirmar que:
\begin{choices}
\choice $MN + MP = 2BM$
\choice $MN + MP = 2CM$
\choice $MN + MP = 2AB$
\choice $MN + MP = 2AD$
\choice $MN + MP = 2AC$
\end{choices}

% figura presente no enunciado
